\subsection{COMMUTATORS IN A LIE GROUP}

PROPOSITION 3. Let G be a finite-dimensional Lie group and \(\mathbf{H}_1\) and \(\mathbf{H}_2\) subgroups of G. Let \(\mathfrak{h}_1, \mathfrak{h}_2\) and \(\mathfrak{h}\) be the Lie subalgebras tangent at \(e\) to \(\mathbf{H}_1\mathbf{H}_2\) and \((\mathbf{H}_1,\mathbf{H}_2)\) respectively. Then \([\mathfrak{h}_1,\mathfrak{h}_2] \subset \mathfrak{h}\) .

Let \(a \in \mathfrak{h}_1\) , \(b \in \mathfrak{h}_2\) . There exist an open neighbourhood I of 0 in K and analytic mappings \(f_1, f_2\) of I into G such that

\[
\begin{array}{c} f _ {1} (0) = f _ {2} (0) = e, \quad f _ {1} (\mathrm{I}) \subset \mathrm{H} _ {1}, \quad f _ {2} (\mathrm{I}) \subset \mathrm{H} _ {2}, \\ (\mathrm{T} _ {0} f _ {1}) 1 = a, \quad (\mathrm{T} _ {0} f _ {2}) 1 = b. \end{array}
\]

We write

\[
f (\lambda , \mu) = (f _ {1} (\lambda), f _ {2} (\mu)) \in (H _ {1}, H _ {2}) \quad \text { for   } \lambda , \mu \text {   in   I. }
\]

We identify an open neighbourhood of \(e\) in \(G\) with an open subset of \(K^r\) by means of a chart which maps \(e\) to 0. Then \(L(G)\) is identified with \(K^r\) . By §5, no. 2, Proposition 1, the expansion of \(f(\lambda, \mu)\) as an integral series about the origin is

\[
f (\lambda , \mu) = \lambda \mu [ a, b ] + \sum_ {i \geqslant 1, j \geqslant 1, i + j \geqslant 2} \lambda^ {i} \mu^ {j} a _ {i j}
\]

where \(a_{ij} \in \mathbf{K}^r\) (the terms in \(\lambda^i\) and \(\mu^j\) in the expansion of \(f(\lambda, \mu)\) are zero because \(f(\lambda, 0) = f(0, \mu) = 0\) ). We fix \(\mu\) in I. Letting \(\lambda\) tend to 0, we see that

\[
\mu [ a, b ] + \sum_ {j \geqslant 2} \mu^ {j} a _ {1 j} \in \mathfrak {h}.
\]

Since this is true for all \(\mu\in\mathbf{I}\) , it follows that \([a,b]\in\mathfrak{h}\) .

Remark. Even if \(H_{1}\) and \(H_{2}\) are connected Lie subgroups of G, the Lie subalgebra of L(G) generated by \([h_{1}, h_{2}]\) is in general distinct from h.

PROPOSITION 4. Let \(\mathbf{G}\) be a finite-dimensional real or complex Lie group. Let \(\mathbf{A}, \mathbf{B}, \mathbf{C}\) be integral subgroups of \(\mathbf{G}\) such that \([\mathrm{L}(\mathrm{A}), \mathrm{L}(\mathrm{C})] \subset \mathrm{L}(\mathrm{C})\) and

\[
[ \mathrm{L} (\mathrm{B}), \mathrm{L} (\mathrm{C}) ] \subset \mathrm{L} (\mathrm{C}).
\]

If \([\mathrm{L}(\mathrm{A}),\mathrm{L}(\mathrm{B})]\subset \mathrm{L}(\mathrm{C})\) , then \((\mathrm{A},\mathrm{B})\subset \mathrm{C}\) . If \([\mathrm{L}(\mathrm{A}),\mathrm{L}(\mathrm{B})] = \mathrm{L}(\mathrm{C})\) , then \((\mathrm{A},\mathrm{B}) = \mathrm{C}\).

Suppose that \([\mathrm{L}(\mathrm{A}),\mathrm{L}(\mathrm{B})]\subset \mathrm{L}(\mathrm{C})\) . The sum \(\mathrm{L}(\mathrm{A}) + \mathrm{L}(\mathrm{B}) + \mathrm{L}(\mathrm{C})\) is a Lie subalgebra of \(\mathrm{L}(\mathrm{G})\) . By considering an integral subgroup of \(\mathbf{G}\) with Lie algebra \(\mathrm{L}(\mathrm{A}) + \mathrm{L}(\mathrm{B}) + \mathrm{L}(\mathrm{C})\) , the problem is reduced to the case where

\[
\mathrm{L} (\mathrm{A}) + \mathrm{L} (\mathrm{B}) + \mathrm{L} (\mathrm{C}) = \mathrm{L} (\mathrm{G})
\]

and G is connected. Then L(C) is an ideal of L(G). Suppose first that G is simply connected. Then C is a normal Lie subgroup of G ( \(\S\) 6, no. 6, Proposition 14). Let \(\phi\) be the canonical morphism of G onto G/C. Then

\[
[ \mathrm{L} (\phi) (\mathrm{L} (\mathrm{A})), \mathrm{L} (\phi) (\mathrm{L} (\mathrm{B})) ] = \{0 \}
\]

and hence \(\phi (\mathbf{A})\) and \(\phi (\mathbf{B})\) commute by the Hausdorff formula; therefore \((\mathrm{A},\mathrm{B})\subset \mathrm{C}\) . In the general case, let \(\mathbf{G}'\) be the universal covering of G and \(\mathbf{A}',\mathbf{B}',\mathbf{C}'\) the integral subgroups of \(\mathbf{G}'\) such that \(\mathrm{L}(\mathrm{A}') = \mathrm{L}(\mathrm{A}),\mathrm{L}(\mathrm{B}') = \mathrm{L}(\mathrm{B}),\) \(\mathrm{L}(\mathrm{C}') = \mathrm{L}(\mathrm{C})\) . Then \((\mathrm{A}',\mathrm{B}')\subset \mathrm{C}'\) and A, B, C are the canonical images of A', B', C' in G, whence \((\mathrm{A},\mathrm{B})\subset \mathrm{C}\) . On the other hand, (A, B) is the underlying set of an integral subgroup of G (§ 6, no. 2, Corollary to Proposition 4) and its Lie algebra contains [L(A), L(B)] (Proposition 3). If

\[
[ \mathrm{L} (\mathrm{A}), \mathrm{L} (\mathrm{B}) ] = \mathrm{L} (\mathrm{C}),
\]

then (A, B) ⊃ C, whence (A, B) = C.

COROLLARY. Let G be a finite-dimensional connected real or complex Lie group with Lie algebra g. The subgroups D \(^{i}\) G (resp. C \(^{i}\) G) are integral subgroups with Lie algebras D \(^{i}\) g (resp. C \(^{i}\) g). If G is simply connected, they are Lie subgroups.

The first assertion follows from Proposition 4 by induction on i. The second follows from the first and § 6, no. 6, Proposition 14.

PROPOSITION 5. Let G be a finite-dimensional real or complex Lie group and A an integral subgroup of G. Then \(\overline{\mathrm{D}\mathrm{A}} = \mathrm{D}\overline{\mathrm{A}}\) . In particular, A is a normal subgroup of \(\overline{\mathbf{A}}\) and \(\overline{\mathbf{A}} / \mathbf{A}\) is commutative.

Let \(\mathfrak{a} = \mathrm{L}(\mathbf{A})\) . Let \(G_{1}\) be the set of \(g \in G\) such that

\[
(\operatorname{Ad} g) x \equiv x (\mathrm{mod} \mathscr {D} \mathfrak {a}) \quad \text { for   all } x \in \mathfrak {a}.
\]

Then \(G_1\) is a closed subgroup of \(G\) . If \(y \in \mathfrak{a}\) , then \(\exp y \in G_1\) , by § 6, no. 4, Corollary 3 (ii) to Proposition 10. Hence \(G_1\) contains \(A\) and therefore \(\bar{A}\) . Thus, for \(g \in \bar{A}\) , \(L(\text{Int } g)\) leaves \(a\) stable and therefore \(\text{Int } g\) leaves \(A\) stable; more precisely, \(L(\text{Int } g)\) defines the identity automorphism of \(a / \mathcal{D}a\) and hence \(\text{Int } g\) defines the identity automorphism of \(A / DA\) . This proves that \((\bar{A}, A) \subset DA\) . With the real Lie group structure on \(G\) , \(\bar{A}\) is a Lie subgroup ( \(\S 8\) , no. 2, Theorem 2); let \(b\) be its Lie algebra. Let \(G_2\) be the set of \(g \in G\) such that

\[
(\operatorname{Ad} g) x \equiv x (\mathrm{mod} \mathscr {D} \mathfrak {a}) \quad \text { for   all } x \in \mathfrak {b}.
\]

By the above, \(\mathbf{G}_2\supset \mathbf{A}\) and hence \(\mathbf{G}_2\supset \overline{\mathbf{A}}\) . Therefore, for \(g\in \overline{\mathbf{A}}\) , Int \(g\) leaves \(\mathrm{D}\overline{\mathbf{A}}\) stable and defines the identity automorphism of \(\overline{\mathbf{A}} / \mathrm{D}\overline{\mathbf{A}}\) . Hence \(\mathrm{D}\overline{\mathbf{A}}\supset \overline{\mathrm{D}\mathbf{A}}\) .

PROPOSITION 6. Suppose that K is ultrametric. Let G be a finite-dimensional Lie group. Let A, B, C be Lie subgroups of G such that \([L(A), L(C)] \subset L(C)\) , \([L(B), L(C)] \subset L(C)\) . If \([L(A), L(B)] \subset L(C)\) , there exist open subgroups \(A'\) , \(B'\) of A, B such that \((A', B') \subset C\) . If \([L(A), L(B)] \subset L(C)\) , there exist open subgroups \(A'\) , \(B'\) , \(C'\) of A, B, C such that \((A', B') = C'\) .

Suppose that \([\mathrm{L}(\mathrm{A}),\mathrm{L}(\mathrm{B})]\subset\mathrm{L}(\mathrm{C})\) . As in the proof of Proposition 4, the problem reduces to the case where \(\mathrm{L}(\mathrm{C})\) is an ideal of \(\mathrm{L}(\mathrm{G})\) . Then, by replacing G by an open subgroup, it reduces to the case where C is normal in G ( \(\S\) 7, no. 1, Proposition 2). Let \(\phi\) be the canonical morphism of G onto G/C. Then

\[
[ \mathrm{L} (\phi) (\mathrm{L} (\mathrm{A})), \mathrm{L} (\phi) (\mathrm{L} (\mathrm{B})) ] = \{0 \}.
\]

By the Hausdorff formula, there exist open subgroups \(\mathbf{A}'\) , \(\mathbf{B}'\) of \(\mathbf{A}\) , \(\mathbf{B}\) such that \(\phi(\mathbf{A}')\) and \(\phi(\mathbf{B}')\) commute, whence \((\mathbf{A}', \mathbf{B}') \subset \mathbf{C}\) . Suppose further that

\[
[ \mathrm{L} (\mathrm{A}), \mathrm{L} (\mathrm{B}) ] = \mathrm{L} (\mathrm{C}).
\]

By Proposition 3, the Lie subalgebra tangent to \((\mathbf{A}',\mathbf{B}')\) at \(e\) contains \(\mathrm{L}(\mathrm{C})\) . Hence \((\mathbf{A}',\mathbf{B}')\) contains a Lie subgroup germ of \(\mathbf{G}\) with Lie algebra \(\mathrm{L}(\mathrm{C})\) . Therefore \((\mathbf{A}',\mathbf{B}')\) is an open subgroup of \(\mathbf{C}\) .

COROLLARY. Suppose that K is ultrametric. Let G be a finite-dimensional Lie group with Lie algebra g. There exists an open subgroup \(G_{0}\) of G such that, for all i, \(D^{i}G_{0}\) (resp. \(C^{i}G_{0}\) ) is a Lie subgroup of G with Lie algebra \(D^{i}g\) (resp. \(C^{i}g\) ).

\begin{enumerate}
\def\labelenumi{(\alph{enumi})}
\item
  By Proposition 3 applied inductively, for every open subgroup \(G_{1}\) of G and for all i, \(D^{i}G_{1}\) contains a Lie subgroup germ of G with Lie algebra \(D^{i}g\) .
\item
  Let \(G'\) be an open subgroup of \(G\) such that, for \(i \leqslant n\) , \(D^i G'\) is a Lie subgroup of \(G\) with Lie algebra \(D^i\mathfrak{g}\) . By Proposition 6, there exist open subgroups \(H_1, H_2\) of \(D^n G'\) , such that \((H_1, H_2)\) is a Lie subgroup with Lie algebra \(D^{n+1}\mathfrak{g}\) . Let \(G''\) be an open subgroup of \(G'\) small enough for \(D^n G'' \subset H_1 \cap H_2\) . Then \(D^{n+1}G'' \subset (H_1, H_2)\) . The relations
\end{enumerate}

\[
\mathrm{D} ^ {0} \mathrm{G} ^ {\prime \prime} \subset \mathrm{D} ^ {0} \mathrm{G} ^ {\prime}, \mathrm{D} ^ {1} \mathrm{G} ^ {\prime \prime} \subset \mathrm{D} ^ {1} \mathrm{G} ^ {\prime}, \dots , \mathrm{D} ^ {n} \mathrm{G} ^ {\prime \prime} \subset \mathrm{D} ^ {n} \mathrm{G} ^ {\prime}, \mathrm{D} ^ {n + 1} \mathrm{G} ^ {\prime \prime} \subset (\mathrm{H} _ {1}, \mathrm{H} _ {2})
\]

prove, using (a), that \(\mathbf{D}^i\mathbf{G}''\) is, for \(i\leqslant n + 1\) , a Lie subgroup of \(\mathbf{G}\) with Lie algebra \(\mathcal{D}^i\mathfrak{g}\) .

\begin{enumerate}
\def\labelenumi{(\alph{enumi})}
\setcounter{enumi}{2}
\item
  There exists an integer p such that \(D^{p}g = D^{p+1}g = \cdots\) . By the above, there exists an open subgroup \(G_{0}\) of G such that \(D^{i}G_{0}\) is, for \(i \leqslant p\) , a Lie subgroup of G with Lie algebra \(D^{i}\mathfrak{g}\) . But, by (a), the same assertion remains true for \(i > p\) since \(D^{p}G_{0} \supset D^{i}G_{0}\) for \(i > p\).
\item
  The argument is similar for the \(C^{i}\) .
\end{enumerate}

